\documentclass[12pt,a4paper]{article}

% --- Преамбула ---
\usepackage[utf8]{inputenc}
\usepackage[T1]{fontenc}
\usepackage[russian]{babel}
\usepackage{amsmath,amssymb}
\usepackage{hyperref}
\usepackage[margin=2.5cm]{geometry}
\usepackage{enumitem}
\usepackage{setspace}
\usepackage{booktabs}
\onehalfspacing

\title{\textbf{Резонансное высвобождение энергии из хвостов собственных мод электрона}\\
\large Физический принцип и экспериментальный протокол}
\author{}
\date{}

\begin{document}
\maketitle

\begin{abstract}
В рамках Субквантовой теории предсказывается, что собственная мода электрона содержит запасённую энергию в виде хвостов causal history. При воздействии резонансного поля, инвертирующего фазу causal history, эта энергия может быть частично высвобождена. Теоретическая оценка даёт коэффициент усиления по энергии около $4$ ($400\%$ от затраченной). Предлагается эксперимент по проверке этого предсказания с использованием одного электрона в вакуумной камере с перестраиваемым электромагнитным полем.
\end{abstract}

\section{Теоретическая основа}
В Субквантовой теории каждая частица обладает собственной модой $\psi_A(k)$ (causal history), определённой на дискретной оси $k = 0, 1, \dots, d_{AA}$. Поперечный спектр $F_{A \to A}(k)$ содержит записи о всех взаимодействиях частицы за всю её историю.

Полная энергия, запасённая в causal history электрона, складывается из двух компонент:
\begin{itemize}
    \item \textbf{Инертная компонента:} энергия, запертая вблизи $k \approx 0$ (отвечает за массу покоя). Она недоступна для высвобождения без разрушения частицы.
    \item \textbf{Гравитационная компонента (хвосты):} энергия, распределённая по $k > K_0$ (за пределами двух слабых нулей). Она доступна для высвобождения при определённых условиях.
\end{itemize}

Для электрона расстояние между двумя слабыми нулями оценивается как $K_0 \approx 0.51$ (в безразмерных единицах). Энергия хвостов составляет примерно $\alpha \cdot m_e c^2$ на одну запись, где $\alpha \approx 1/137$ --- постоянная тонкой структуры.

\section{Механизм высвобождения}
При нормальном ходе времени causal history растёт, и энергия запасается. Градиентный спуск минимизирует энергию сети:
\begin{equation}
    F(k, t+1) = F(k, t) - \alpha \cdot \frac{\partial E}{\partial F(k)}.
\end{equation}

Если приложить внешнее поле, инвертирующее фазу градиентного спуска ($F \to F + \alpha \cdot \partial E/\partial F$), система начинает двигаться в сторону увеличения энергии. Хвосты causal history, ранее запасённые, начинают «разворачиваться», высвобождая энергию.

Условие резонанса: частота внешнего поля должна совпадать с собственной частотой causal history электрона:
\begin{equation}
    \nu_{\text{res}} = \frac{c}{2 \cdot K_0 \cdot \delta},
\end{equation}
где $\delta$ --- фундаментальная длина (порядка планковской). Точное значение $\nu_{\text{res}}$ подлежит экспериментальному определению.

\section{Теоретическая оценка выхлопа}
Полная энергия causal history одного электрона:
\begin{equation}
    E_{\text{total}} \approx N_{\text{records}} \cdot (\alpha \cdot m_e c^2),
\end{equation}
где $N_{\text{records}}$ --- число записей в causal history (порядка $10^{17}$ для электрона, существовавшего значительную часть возраста Вселенной).

Численная оценка:
\begin{equation}
    E_{\text{total}} \approx 4.3 \times 10^{17} \times (0.511\,\text{МэВ} / 137) \approx 1.6 \times 10^{21}\,\text{эВ} \approx 260\,\text{Дж}.
\end{equation}

Энергия, доступная для высвобождения (хвосты), оценивается как $\sim 260$~Дж.

Энергия, необходимая для инверсии фазы causal history (затраты), оценивается как $E_{\text{total}} / 4 \approx 65$~Дж, исходя из модели, где КПД процесса составляет $25\%$ (коэффициент усиления $= 4$).

\section{Экспериментальный протокол}

\subsection{Цель}
Измерить энергию, высвобождаемую одним электроном при резонансном воздействии поля, инвертирующего фазу causal history.

\subsection{Установка}
\begin{itemize}
    \item Электронная пушка (источник одиночных электронов с известной кинетической энергией $\sim 10$~кэВ).
    \item Вакуумная камера с системой электродов, создающих переменное электрическое поле с перестраиваемой частотой (диапазон: от $1$~МГц до $10$~ГГц, с возможностью расширения).
    \item Калориметр для измерения энергии электрона до и после взаимодействия с полем.
    \item Система синхронизации: электрон должен пролетать через область поля в определённой фазе.
\end{itemize}

\subsection{Процедура}
\begin{enumerate}[label=\arabic*.]
    \item Калибровка: измеряется энергия электрона без поля (контроль).
    \item Сканирование частоты: частота поля меняется в заданном диапазоне. Для каждой частоты измеряется энергия электрона после взаимодействия.
    \item Регистрация резонанса: при совпадении частоты поля с $\nu_{\text{res}}$ энергия электрона после взаимодействия должна возрасти (выхлоп).
    \item Измерение коэффициента усиления: отношение выхлоп / затраты $= (E_{\text{после}} - E_{\text{до}}) / E_{\text{поля}}$, где $E_{\text{поля}}$ --- энергия, затраченная на создание поля в объёме взаимодействия.
\end{enumerate}

\subsection{Ожидаемый результат}
\begin{itemize}
    \item При частотах, далёких от $\nu_{\text{res}}$: $E_{\text{после}} \approx E_{\text{до}}$ (выхлоп отсутствует).
    \item При $\nu \approx \nu_{\text{res}}$: $E_{\text{после}} > E_{\text{до}}$. Коэффициент усиления достигает $\sim 4$.
    \item Форма резонансной кривой: пик с шириной, определяемой добротностью causal history электрона.
\end{itemize}

\section{Обсуждение}
Если предсказанный эффект будет обнаружен, это станет первым прямым экспериментальным подтверждением Субквантовой теории и Теории Мод. Это также откроет путь к созданию принципиально новых источников энергии, основанных на управлении causal history.

Эксперимент с одним электроном абсолютно безопасен и не требует крупных энергетических затрат. Он может быть выполнен в стандартной лаборатории ядерной физики или физики элементарных частиц.

\section{Заключение}
Предложен протокол эксперимента по резонансному высвобождению энергии из хвостов собственных мод электрона. Теоретическая оценка предсказывает коэффициент усиления $\sim 4$. Экспериментальная проверка этого предсказания является приоритетной задачей для верификации Субквантовой теории.

\end{document}